% ============================================
% LUC-05
% Gerado de molambudos.md — build_miolo.py
% ============================================
\section*{\textbf{LUC-05} --- O Sonho de Lúcia}
\markboth{LUC-05 --- O Sonho de Lúcia}{}

\noindent \textit{A contaminação da investigadora}

\noindent \textbf{Abril de 2026 --- Belo Horizonte}

\vspace{0.5em}\noindent\rule{\textwidth}{0.5pt}\vspace{0.5em}

\noindent Ela sonhou com a vala.

\noindent Não como quem vê um filme. Como quem \textit{está} lá. Os pés descalços na terra fria. O cheiro adocicado que entra pelos poros. O céu escuro, sem estrelas, sem lua, como se o mundo tivesse esquecido de acender as luzes.

\noindent Ela ouviu a mastigação antes de ver a criatura.

\noindent O som de osso sendo quebrado. De carne sendo triturada. Um som molhado, paciente, rítmico --- como quem come devagar porque não tem pressa, porque a comida não vai fugir.

\noindent Ela andou até a borda da vala.

\noindent Lá embaixo, a pilha de corpos se moveu.

\noindent A criatura ergueu a cabeça.

\noindent O rosto quase humano. Pele cinzenta esticada sobre os ossos. A boca suja de sangue. Os olhos amarelos brilhando no escuro.

\noindent Ela sorriu.

\noindent \textit{\textquotedblleft{}Finalmente\textquotedblright{}\textit{, disse a criatura. }\textquotedblleft{}Estava esperando a senhora.\textquotedblright{}}

\noindent \textit{\textquotedblleft{}Você sabia que eu ia vir?\textquotedblright{}}

\noindent \textit{\textquotedblleft{}Sempre soube. Desde a primeira página.\textquotedblright{}}

\noindent A criatura subiu. O corpo imenso, a gaiola de costelas expostas, o órgão pulsando entre os ossos. Ela parou na frente de Lúcia.

\noindent \textit{\textquotedblleft{}A senhora leu o diário.\textquotedblright{}}

\noindent \textit{\textquotedblleft{}...Sim.\textquotedblright{}}

\noindent \textit{\textquotedblleft{}O que achou?\textquotedblright{}}

\noindent \textit{\textquotedblleft{}É a história mais triste que já li.\textquotedblright{}}

\noindent A criatura inclinou a cabeça.

\noindent \textit{\textquotedblleft{}Não é triste. É verdade. A senhora confunde as duas coisas.\textquotedblright{}}

\noindent \textit{\textquotedblleft{}O que você quer de mim?\textquotedblright{}}

\noindent \textit{\textquotedblleft{}Nada. Já tenho o que queria.\textquotedblright{}}

\noindent \textit{\textquotedblleft{}O quê?\textquotedblright{}}

\noindent \textit{\textquotedblleft{}A senhora.\textquotedblright{}}

\noindent A criatura tocou o peito de Lúcia com um dedo longo. O dedo atravessou a pele, os ossos. Não doeu. Apenas... entrou.

\noindent \textit{\textquotedblleft{}Agora a senhora é minha.\textquotedblright{}}

\noindent \textit{\textquotedblleft{}Não\textquotedblright{}\textit{, disse Lúcia. }\textquotedblleft{}Não sou.\textquotedblright{}}

\noindent \textit{\textquotedblleft{}Já é. Desde o momento em que abriu o diário. A senhora só não aceitou ainda.\textquotedblright{}}

\noindent Lúcia acordou.

\noindent O quarto estava escuro. O relógio marcava 3:14.

\noindent A doçura podre estava no travesseiro.

\noindent Ela sentou na cama, ofegante. Olhou para o canto do quarto. A sombra parecia mais densa.

\noindent Ela acendeu a luz.

\noindent Nada.

\noindent Mas o cheiro não foi embora.
\noindent Ela sentou na cama, ofegante. Olhou para o canto do quarto. A sombra parecia mais densa. E havia um som. Um som que ela conhecia do relatório de Joaquim — um canto fino, agudo, como uma mulher chorando muito longe. O canto de uma coruja. A Rasga Mortalha. Ela não viu o pássaro. Mas ouviu. Ouviu as asas se abrindo, o bater no ar parado do quarto.

\noindent Ela ficou sentada na cama, com a luz acesa, até o sol nascer.

\vspace{0.5em}\noindent\rule{\textwidth}{0.5pt}\vspace{0.5em}

\noindent Na manhã seguinte, ela foi ao oftalmologista.

\noindent \textit{\textquotedblleft{}Seu olho direito está com uma leve descoloração\textquotedblright{}\textit{, disse o médico. }\textquotedblleft{}Nada grave. Apenas uma alteração na pigmentação da íris. Já nasceu assim?\textquotedblright{}}

\noindent \textit{\textquotedblleft{}Não.\textquotedblright{}}

\noindent \textit{\textquotedblleft{}Hmm. Pode ser medicamento. A senhora toma algum remédio?\textquotedblright{}}

\noindent \textit{\textquotedblleft{}Não.\textquotedblright{}}

\noindent \textquotedblleft{\textit{Estranho.\textquotedblright{}} Ele anotou algo no prontuário. }\textquotedblleft{\textit{Vou pedir uns exames. Só para garantir.\textquotedblright{}}}

\noindent Lúcia olhou no espelho do consultório.

\noindent O olho direito estava mais claro.
Amarelado.

\noindent Ela não disse nada.

\noindent Mas, ao sair, passou a mão no rosto e sentiu a pele mais fria do que deveria.
\noindent Alguma coisa havia assinado por ela.

\noindent O ciclo tinha começado.

\vspace{0.5em}
\noindent\textit{\textbf{↪ Rotas:} \rota{LUC-Contraprova} \rotasep \rota{CONT-04}}